\documentclass[11pt, letterpaper]{article}

\usepackage{amsmath, amssymb}
\usepackage{booktabs}
\usepackage{geometry}
\usepackage{hyperref}
\usepackage{xcolor}
\usepackage{enumitem}
\usepackage{tcolorbox}
\usepackage{graphicx}
\usepackage{subcaption}
\usepackage{float}
\usepackage{array}

\geometry{margin=1in}

\tcbuselibrary{skins, breakable}
\newtcolorbox{agentbox}[2][]{
  title={#2}, colback=blue!4!white, colframe=blue!50!black,
  fonttitle=\bfseries, breakable, #1
}
\newtcolorbox{answerbox}[2][]{
  title={#2}, colback=green!4!white, colframe=green!50!black,
  fonttitle=\bfseries, breakable, #1
}
\newtcolorbox{findingbox}[2][]{
  title={#2}, colback=yellow!6!white, colframe=orange!60!black,
  fonttitle=\bfseries, breakable, #1
}

\newcommand{\Dt}{\Delta t}
\newcommand{\lrt}{\lambda^{RT}}
\newcommand{\pda}{p^{DA}}
\newcommand{\figdir}{../results/figures/2024-07-08_5d}

\title{\textbf{Agent Optimization Formulations \& Case Study Results} \\[4pt]
\large Iteration-Free Two-Timescale Market Participation via Explicit mp-GNE}
\author{Autonomous Market RHG}
\date{\today}

\begin{document}
\maketitle
\tableofcontents
\newpage

%% ============================================================
\section{Problem Overview}
%% ============================================================

\begin{answerbox}{Setup}

A microgrid behind a single Point of Common Coupling (PCC) participates in a two-settlement wholesale market (Day-Ahead and Real-Time). It contains $N = 4$ heterogeneous agents that must agree on a joint power profile at every 5-minute interval over a receding horizon of $H = 6$ steps (30 minutes).

\textbf{Agents:}
\begin{itemize}[nosep]
    \item Agent 1: Vanadium Redox Flow Battery (VRFB), 1 MW / 7 MWh
    \item Agent 2: PV + Li-ion battery, 1 MW PV + 1 MW / 4 MWh storage
    \item Agent 3: PEM electrolyzer, 2 MW, fast ramp
    \item Agent 4: Alkaline electrolyzer, 1.5 MW, slow ramp
\end{itemize}

\textbf{Market: PJM-RTO}, week of 2024-07-08, $L_{\max} = 1500$ kW PCC import cap.

\textbf{Two timescales:}
\begin{itemize}[nosep]
    \item \emph{Day-Ahead} (hourly, $H_{DA} = 24$): centralized CVXPY clears the daily schedule.
    \item \emph{Real-Time} (5-minute, $H = 6$): distributed Generalized Nash Equilibrium solved by FACET-mpGNE (iteration-free explicit map) with ADMM as a baseline.
\end{itemize}

\end{answerbox}

%% ============================================================
\section{Global Constraints}
%% ============================================================

\begin{answerbox}{Shared Feasibility: Global Coupling at the PCC}

The defining feature of the Generalized Nash Equilibrium (GNE) is the shared PCC. The sum of all agent net power must not exceed the transformer rating:
\[
  \sum_{i=1}^N p_{i,k} \;\le\; L_{\max}, \quad \forall k \in \{1, \ldots, H\}
\]
where $p_{i,k}$ is agent $i$'s net power at horizon step $k$ (positive = import, negative = export), and $L_{\max}$ is the PCC import limit (kW).

\textbf{Note on $L_{\min}$ (export cap):} Including a lower bound $\sum_i p_{i,k} \ge L_{\min}$ would double the coupling constraints and cause a combinatorial explosion of Critical Regions (CRs) in the explicit mpQP. Since each agent's per-unit power bounds already cap the achievable export, we drop $L_{\min}$ entirely. The simulation confirms this is non-binding.

\end{answerbox}

%% ============================================================
\section{Agent Optimization Formulations (Full Detail)}
%% ============================================================

Each agent solves its own quadratic program every 5 minutes.  Each agent sees only
(i) its own physical state, (ii) its own DA schedule from the previous day's market
clear, (iii) the shared RTM LMP forecast, and (iv) the \emph{aggregate power of the
other three agents} as an exogenous input that enters only through the PCC coupling
constraint.  No agent observes the other agents' objectives, efficiencies, or
inventory levels.

The shared uncertain inputs at each RT step are:
\begin{itemize}[nosep]
  \item $\lambda^{RT}_k$, $k=1,\ldots,6$: RTM LMP forecast over the 30-min horizon,
        corrupted by $\mathcal{N}(0,\,20^2)$\,\$/MWh noise.
  \item $p^{DA}_{i,k}$: each agent's own DA schedule (6 values, contracted at DAM clear, exact).
  \item Own initial state $\text{SoC}_{i,0}$ or $\text{Inv}_{i,0}$: plant measurement, exact.
\end{itemize}

The coupling constraint from agent $i$'s perspective is:
\[
  p_{i,k} \;\le\; L_{\max} - \textstyle\sum_{j \ne i} p_{j,k}, \quad k = 1,\ldots,6,
\]
where $L_{\max} = 1500$\,kW and $\sum_{j \ne i} p_{j,k}$ is the aggregate of the other three agents' decisions (treated as a given at each iteration).

\textbf{DA schedule as a scalar parameter.}
The DAM clears once per day and produces 24 hourly values $\{p^{DA}_{i,h}\}_{h=0}^{23}$ per agent.  Within any 30-minute RT horizon ($H=6$ steps), all steps fall in the same hourly bucket, so the DA reference is a single scalar
\[
  p^{DA}_{i} \;=\; p^{DA}_{i,\,\lfloor t/12 \rfloor}
\]
updated once per hour.  Each agent's private parameter vector is therefore:
\[
  \theta_i \;=\; \bigl[\textstyle\sum_{j\ne i}p_{j,k}\bigr]_{k=1}^{6}
              \oplus\; \text{state}_{i,0}
              \oplus\; [\lambda^{RT}_k]_{k=1}^{6}
              \oplus\; p^{DA}_{i}
              \;\in\; \mathbb{R}^{14},
\]
and for Agent 2 (PV+Batt), which additionally carries 6 PV-forecast values: $\theta_2 \in \mathbb{R}^{20}$.  Reducing DA from $H=6$ values to 1 scalar cuts $|\theta_i|$ from 19 to 14 (and 25 to 20), dramatically reducing the Critical Region count while keeping the mpQP rebuilt only \emph{once per day}.

%% ------------------------------------------------------------
\subsection{Agent 1: VRFB (Vanadium Redox Flow Battery)}
%% ------------------------------------------------------------

\begin{agentbox}{Agent 1 — Bidirectional flow battery, 1\,MW / 7\,MWh}

\textbf{Fixed parameters (never change during the simulation):}
\begin{center}
\small
\begin{tabular}{lll}
\toprule
Symbol & Value & Meaning \\
\midrule
$P^{\max}$ & 1000\,kW & Rated charge and discharge power \\
$E_{\min}$ & 700\,kWh & SoC lower bound (10\% of 7\,MWh) \\
$E_{\max}$ & 6300\,kWh & SoC upper bound (90\% of 7\,MWh) \\
$\eta$     & 0.75      & Charge/discharge efficiency (linear model) \\
$\Dt$      & 5/60\,h   & Step duration \\
$\gamma$   & 0.1\,\$/kW$^2$ & DA-tracking imbalance penalty \\
\bottomrule
\end{tabular}
\end{center}

\textbf{Uncertain inputs (updated every RT step):}
\begin{itemize}[nosep]
  \item $\text{SoC}_0$\,: measured state of charge at the start of the step (initial value: 3500\,kWh).
  \item $\lambda^{RT}_k$, $k=1\ldots6$: RTM LMP forecast, $\sigma = 20$\,\$/MWh.
  \item $p^{DA}_k$, $k=1\ldots6$: own DA schedule from DAM clear (kW).
\end{itemize}

\textbf{Decision variable:} $p_k \in \mathbb{R}$, $k=1,\ldots,6$ (kW; positive = charging from grid, negative = discharging to grid).

\textbf{State dynamics:}
\[
\text{SoC}_k \;=\; \text{SoC}_0 \;+\; \eta\,\Dt\,\sum_{j=1}^{k} p_j, \qquad k = 1,\ldots,6.
\]
(Self-discharge neglected: $\alpha = 1$.)

\textbf{Objective — minimise energy cost plus deviation from DA schedule:}
\[
\min_{\{p_k\}}\;\sum_{k=1}^{6}\Big[\;\underbrace{\Dt\,\lambda^{RT}_k\,p_k}_{\text{energy cost}}
  \;+\; \underbrace{\tfrac{1}{2}\,\gamma\,(p_k - p^{DA}_k)^2}_{\text{DA imbalance penalty}}\;\Big]
\]
The unconstrained optimum is $p^*_k = p^{DA}_k - \Dt\,\lambda^{RT}_k\,/\,\gamma$: when RT prices are high the battery discharges (exports), when low it charges.

\textbf{Constraints:}
\begin{align*}
  &\text{Power limits:}   && -1000 \;\le\; p_k \;\le\; 1000 \quad\text{kW} \\
  &\text{SoC upper:}      && \text{SoC}_0 + \eta\,\Dt\,\textstyle\sum_{j=1}^k p_j \;\le\; 6300 \quad\text{kWh} \\
  &\text{SoC lower:}      && \text{SoC}_0 + \eta\,\Dt\,\textstyle\sum_{j=1}^k p_j \;\ge\; 700  \quad\text{kWh} \\
  &\text{PCC coupling:}   && p_k \;\le\; 1500 - \textstyle\sum_{j\ne 1} p_{j,k} \quad\text{kW}
\end{align*}
for $k = 1,\ldots,6$.\quad Total: 24 local + 6 coupling = \textbf{30 constraints}.

\end{agentbox}

%% ------------------------------------------------------------
\subsection{Agent 2: PV + Li-ion Battery}
%% ------------------------------------------------------------

\begin{agentbox}{Agent 2 — 1\,MW rooftop PV with 1\,MW / 4\,MWh Li-ion storage}

\textbf{Fixed parameters:}
\begin{center}
\small
\begin{tabular}{lll}
\toprule
Symbol & Value & Meaning \\
\midrule
$P^{\max}_{\text{batt}}$ & 1000\,kW & Battery inverter rated power \\
$P^{\max}_{\text{PV}}$   & 1000\,kW & PV nameplate capacity \\
$E_{\min}$ & 400\,kWh  & SoC lower bound (10\% of 4\,MWh) \\
$E_{\max}$ & 3800\,kWh & SoC upper bound (95\% of 4\,MWh) \\
$\eta$     & 0.92      & Battery charge/discharge efficiency \\
$\Dt$      & 5/60\,h   & Step duration \\
$\gamma$   & 0.1\,\$/kW$^2$ & DA-tracking imbalance penalty \\
\bottomrule
\end{tabular}
\end{center}

\textbf{Uncertain inputs (updated every RT step):}
\begin{itemize}[nosep]
  \item $\text{SoC}_0$: measured battery state of charge (initial: 2000\,kWh).
  \item $\lambda^{RT}_k$, $k=1\ldots6$: RTM LMP forecast, $\sigma = 20$\,\$/MWh.
  \item $p^{DA}_k$, $k=1\ldots6$: own DA schedule (kW).
  \item $g^{PV}_k$, $k=1\ldots6$: PV generation forecast (kW).
        Received as a \emph{hourly} update (zero-order hold over 12 steps), with additive $\mathcal{N}(0,\,150^2)$\,kW noise clipped to $[0,\,1000]$\,kW.
        This is the only agent whose problem depends on the PV forecast.
\end{itemize}

\textbf{Decision variable:} $p_k \in \mathbb{R}$ (kW net PCC power; positive = importing from grid, negative = exporting to grid).

\textbf{State dynamics:}
PV generation $g^{PV}_k$ flows into the battery before the net grid power.  The effective battery state is:
\[
\text{SoC}_k \;=\; \text{SoC}_0 \;+\; \eta\,\Dt\,\sum_{j=1}^{k}\bigl(p_j + g^{PV}_j\bigr), \qquad k = 1,\ldots,6.
\]
The decision variable $p_k$ is the \emph{net} import: when $p_k < g^{PV}_k$ the surplus PV charges the battery; when $p_k + g^{PV}_k < 0$ the battery discharges to the grid.

\textbf{Objective — minimise energy bill plus DA deviation (PV is free):}
\[
\min_{\{p_k\}}\;\sum_{k=1}^{6}\Big[\;\Dt\,\lambda^{RT}_k\,p_k
  \;+\; \tfrac{1}{2}\,\gamma\,(p_k - p^{DA}_k)^2\;\Big]
\]

\textbf{Constraints:}
\begin{align*}
  &\text{Battery inverter upper:} && p_k \;\le\; P^{\max}_{\text{batt}} - g^{PV}_k \quad\text{(PV reduces import headroom)} \\
  &\text{Battery inverter lower:} && p_k \;\ge\; -P^{\max}_{\text{batt}} - g^{PV}_k \quad\text{(PV increases export capacity)} \\
  &\text{SoC upper:}   && \text{SoC}_0 + \eta\,\Dt\,\textstyle\sum_{j=1}^k (p_j + g^{PV}_j) \;\le\; 3800 \quad\text{kWh} \\
  &\text{SoC lower:}   && \text{SoC}_0 + \eta\,\Dt\,\textstyle\sum_{j=1}^k (p_j + g^{PV}_j) \;\ge\; 400  \quad\text{kWh} \\
  &\text{PCC coupling:}&& p_k \;\le\; 1500 - \textstyle\sum_{j\ne 2} p_{j,k} \quad\text{kW}
\end{align*}
for $k = 1,\ldots,6$.

The PV forecast $g^{PV}_k$ shifts the power bounds in both directions and also appears on the right-hand side of both SoC constraints — it is the only uncertain input that enters \emph{four} constraint blocks simultaneously.

Total: 24 local + 6 coupling = \textbf{30 constraints}.

\end{agentbox}

%% ------------------------------------------------------------
\subsection{Agent 3: PEM Electrolyzer}
%% ------------------------------------------------------------

\begin{agentbox}{Agent 3 — PEM electrolyzer, 2\,MW, fast-ramping H$_2$ production}

\textbf{Fixed parameters:}
\begin{center}
\small
\begin{tabular}{lll}
\toprule
Symbol & Value & Meaning \\
\midrule
$P^{\max}$ & 2000\,kW & Maximum power draw (always consuming, $p_k \ge 0$) \\
$\eta^{\text{prod}}$ & 0.020\,kg/kWh & H$_2$ production yield (50\,kWh per kg) \\
$M_{\min}$ & 0\,kg    & Tank lower bound \\
$M_{\max}$ & 800\,kg  & Tank upper bound \\
$\pi_{H_2}$ & 4.00\,\$/kg & H$_2$ offtake contract price \\
$q$        & $400/288 \approx 1.389$\,kg/step & Constant offtake per step (400\,kg/day contract) \\
$\Dt$      & 5/60\,h  & Step duration \\
$\gamma$   & 0.1\,\$/kW$^2$ & DA-tracking imbalance penalty \\
\bottomrule
\end{tabular}
\end{center}

\textbf{Uncertain inputs (updated every RT step):}
\begin{itemize}[nosep]
  \item $\text{Inv}_0$: measured H$_2$ tank level (initial: 200\,kg).
  \item $\lambda^{RT}_k$, $k=1\ldots6$: RTM LMP forecast, $\sigma = 20$\,\$/MWh.
  \item $p^{DA}_k$, $k=1\ldots6$: own DA schedule (kW).
\end{itemize}

\textbf{Decision variable:} $p_k \ge 0$, $k=1,\ldots,6$ (kW consumed; always non-negative — PEM only draws power).

\textbf{Inventory dynamics with constant offtake:}
\[
\text{Inv}_k \;=\; \text{Inv}_0 \;+\; \eta^{\text{prod}}\,\Dt\,\sum_{j=1}^{k} p_j \;-\; q\,k, \qquad k = 1,\ldots,6.
\]
The cumulative offtake $q\,k$ drains the tank at a constant rate.  When $p_k$ is large the tank refills; when $p_k$ is small (high LMP) the tank drains.

\textbf{Objective — minimise energy cost minus H$_2$ revenue, plus DA deviation:}
\[
\min_{\{p_k \ge 0\}}\;\sum_{k=1}^{6}\Big[\;\underbrace{\Dt\,\lambda^{RT}_k\,p_k}_{\text{electricity cost}}
  \;-\; \underbrace{\pi_{H_2}\,\eta^{\text{prod}}\,\Dt\,p_k}_{\text{H}_2\text{ revenue}}
  \;+\; \underbrace{\tfrac{1}{2}\,\gamma\,(p_k - p^{DA}_k)^2}_{\text{DA imbalance penalty}}\;\Big]
\]
The H$_2$ revenue term makes the effective energy price $\lambda^{RT}_k - \pi_{H_2}\,\eta^{\text{prod}}\times 1000 = \lambda^{RT}_k - 80$\,\$/MWh.  The PEM runs at full power when $\lambda^{RT}_k < 80$\,\$/MWh and cuts back above that.

\textbf{Constraints:}
\begin{align*}
  &\text{Power bounds:}     && 0 \;\le\; p_k \;\le\; 2000 \quad\text{kW} \\
  &\text{Tank upper:}       && \text{Inv}_0 + \eta^{\text{prod}}\,\Dt\,\textstyle\sum_{j=1}^k p_j - q\,k \;\le\; 800 \quad\text{kg} \\
  &\text{Tank lower:}       && \text{Inv}_0 + \eta^{\text{prod}}\,\Dt\,\textstyle\sum_{j=1}^k p_j - q\,k \;\ge\; 0   \quad\text{kg} \\
  &\text{PCC coupling:}     && p_k \;\le\; 1500 - \textstyle\sum_{j\ne 3} p_{j,k} \quad\text{kW}
\end{align*}
for $k = 1,\ldots,6$.

Total: 24 local + 6 coupling = \textbf{30 constraints}.

\end{agentbox}

%% ------------------------------------------------------------
\subsection{Agent 4: Alkaline Electrolyzer}
%% ------------------------------------------------------------

\begin{agentbox}{Agent 4 — Alkaline electrolyzer, 1.5\,MW, lower-efficiency H$_2$ production}

\textbf{Fixed parameters:}
\begin{center}
\small
\begin{tabular}{lll}
\toprule
Symbol & Value & Meaning \\
\midrule
$P^{\max}$ & 1500\,kW & Maximum power draw (always consuming, $p_k \ge 0$) \\
$\eta^{\text{prod}}$ & 0.018\,kg/kWh & H$_2$ production yield ($\approx$55.6\,kWh per kg) \\
$M_{\min}$ & 0\,kg     & Tank lower bound \\
$M_{\max}$ & 1200\,kg  & Tank upper bound (larger buffer than PEM) \\
$\pi_{H_2}$ & 4.00\,\$/kg & H$_2$ offtake contract price \\
$q$        & $300/288 \approx 1.042$\,kg/step & Constant offtake per step (300\,kg/day contract) \\
$\Dt$      & 5/60\,h   & Step duration \\
$\gamma$   & 0.1\,\$/kW$^2$ & DA-tracking imbalance penalty \\
\bottomrule
\end{tabular}
\end{center}

\textbf{Uncertain inputs (updated every RT step):}
\begin{itemize}[nosep]
  \item $\text{Inv}_0$: measured H$_2$ tank level (initial: 320\,kg).
  \item $\lambda^{RT}_k$, $k=1\ldots6$: RTM LMP forecast, $\sigma = 20$\,\$/MWh.
  \item $p^{DA}_k$, $k=1\ldots6$: own DA schedule (kW).
\end{itemize}

\textbf{Decision variable:} $p_k \ge 0$, $k=1,\ldots,6$ (kW consumed; always non-negative).

\textbf{Inventory dynamics with constant offtake:}
\[
\text{Inv}_k \;=\; \text{Inv}_0 \;+\; \eta^{\text{prod}}\,\Dt\,\sum_{j=1}^{k} p_j \;-\; q\,k, \qquad k = 1,\ldots,6.
\]

\textbf{Objective — minimise energy cost minus H$_2$ revenue, plus DA deviation:}
\[
\min_{\{p_k \ge 0\}}\;\sum_{k=1}^{6}\Big[\;\Dt\,\lambda^{RT}_k\,p_k
  \;-\; \pi_{H_2}\,\eta^{\text{prod}}\,\Dt\,p_k
  \;+\; \tfrac{1}{2}\,\gamma\,(p_k - p^{DA}_k)^2\;\Big]
\]
The effective break-even price is $\pi_{H_2}\,\eta^{\text{prod}}\times 1000 = 72$\,\$/MWh, lower than PEM's 80\,\$/MWh.  Alkaline therefore reduces production sooner as RT prices rise — consistent with its lower energy-to-H$_2$ conversion efficiency.

\textbf{Constraints:}
\begin{align*}
  &\text{Power bounds:}     && 0 \;\le\; p_k \;\le\; 1500 \quad\text{kW} \\
  &\text{Tank upper:}       && \text{Inv}_0 + \eta^{\text{prod}}\,\Dt\,\textstyle\sum_{j=1}^k p_j - q\,k \;\le\; 1200 \quad\text{kg} \\
  &\text{Tank lower:}       && \text{Inv}_0 + \eta^{\text{prod}}\,\Dt\,\textstyle\sum_{j=1}^k p_j - q\,k \;\ge\; 0    \quad\text{kg} \\
  &\text{PCC coupling:}     && p_k \;\le\; 1500 - \textstyle\sum_{j\ne 4} p_{j,k} \quad\text{kW}
\end{align*}
for $k = 1,\ldots,6$.

Total: 24 local + 6 coupling = \textbf{30 constraints}.

\end{agentbox}

%% ── Summary table ─────────────────────────────────────────────────────────────
\subsection{Agent Summary}

\begin{table}[H]
\centering
\caption{All four agents at a glance.  Each agent knows only its own row.}
\label{tab:agent_summary}
\small
\begin{tabular}{l|cccc}
\toprule
 & Agent 1 (VRFB) & Agent 2 (PV+Batt) & Agent 3 (PEM) & Agent 4 (Alk) \\
\midrule
\textbf{Power range (kW)} & $[-1000,\,+1000]$ & $[-1000,\,+1000]^{*}$ & $[0,\,2000]$ & $[0,\,1500]$ \\
\textbf{Energy / Tank}    & [700, 6300]\,kWh & [400, 3800]\,kWh & [0, 800]\,kg & [0, 1200]\,kg \\
\textbf{Init state}       & 3500\,kWh & 2000\,kWh & 200\,kg & 320\,kg \\
\textbf{Efficiency $\eta$} & 0.75 & 0.92 & 0.020\,kg/kWh & 0.018\,kg/kWh \\
\textbf{H$_2$ price}      & — & — & 4.00\,\$/kg & 4.00\,\$/kg \\
\textbf{Daily offtake}    & — & — & 400\,kg/day & 300\,kg/day \\
\textbf{Break-even LMP}   & — & — & 80\,\$/MWh & 72\,\$/MWh \\
\textbf{Uncertain inputs} & $\lambda^{RT}$, $p^{DA}$, SoC$_0$ & $\lambda^{RT}$, $p^{DA}$, SoC$_0$, $g^{PV}$ & $\lambda^{RT}$, $p^{DA}$, Inv$_0$ & $\lambda^{RT}$, $p^{DA}$, Inv$_0$ \\
\textbf{Local constraints} & 24 & 24 & 24 & 24 \\
\bottomrule
\multicolumn{5}{l}{$^*$Effective limits shift with PV: upper $= 1000 - g^{PV}_k$, lower $= -(1000 + g^{PV}_k)$.}
\end{tabular}
\end{table}

\textbf{Note on ramp constraints:} Alkaline electrolyzers have slow ramps ($\pm30$\,kW/min) and PEM is faster ($\pm100$\,kW/min) in practice.  These are intentionally omitted to keep the mpQP small — adding ramp rows would grow the local constraint block from 24 to 36 rows and substantially increase the number of Critical Regions.  The imbalance penalty $\tfrac{1}{2}\gamma(p_k - p^{DA}_k)^2$ already penalises large step-changes from the DA schedule.

\newpage

%% ============================================================
\section{Two-Timescale Architecture}
%% ============================================================

\begin{answerbox}{Day-Ahead (hourly) $\to$ Real-Time (5-min) Pipeline}

The market participation problem has two natural timescales which we treat separately:

\textbf{Layer 1 — Day-Ahead Market clearing} (hourly, $H_{DA} = 24$, $\Dt = 1$ h):
\begin{enumerate}[nosep]
    \item Each agent's QP is rebuilt with the hourly cadence: $\gamma_{\text{imb}} = 10^{-4}$ (weak DA reference; this is the DAM solve itself).
    \item All four agents' decisions are co-optimized as a single centralized QP (CVXPY + CLARABEL) using:
    \begin{itemize}[nosep]
        \item hourly DAM LMP $\lambda^{DA}_h$, $h = 1\ldots24$,
        \item hourly PV forecast $g^{PV,DA}_h$,
        \item current end-of-day states $(\text{SoC}_v^{\text{end}}, \text{SoC}_b^{\text{end}}, \text{Inv}_{PEM}^{\text{end}}, \text{Inv}_{Alk}^{\text{end}})$.
    \end{itemize}
    Each agent's DA schedule is determined entirely by the centralized optimizer — no bidding-range constraints are imposed; the optimizer selects the cost-minimizing schedule subject only to the agents' physical limits and the PCC cap.
    \item Output: 24 hourly DA quantities $\{p^{DA}_{i,h}\}_{h=1}^{24}$ per agent.
\end{enumerate}

\textbf{Layer 2 — Daily mpQP rebuild with tight DA box}:

After the DAM clears, for each agent $i$ we extract the tight DA box from the actual day's schedule:
\[
\underline{d}_i \;=\; \min_h p^{DA}_{i,h} - 1\text{ kW}, \qquad
\overline{d}_i \;=\; \max_h p^{DA}_{i,h} + 1\text{ kW}
\]
This box is dramatically smaller than the agent's static DA-bidding box (e.g.\ Agent 1 VRFB collapses to $[-21, -19]$ kW, only 2 kW wide). The explicit mpQP is then re-solved using this per-day, per-agent tight box. The resulting Critical Regions are valid for any hourly DA value within the day's actual schedule.

\textbf{Layer 3 — Real-Time GNE} (5-min, $H = 6$, $\Dt = 5/60$ h):

At each RT step $t$:
\begin{enumerate}[nosep]
    \item Look up the current hourly DA value: $p^{DA}_{i,k} = p^{DA}_{i,\,\lfloor(t+k)/12\rfloor}$ for each agent and horizon step.
    \item Pack $\theta_i$ (own state, RT LMP forecast, current DA hourly value, PV forecast for Agent 2) and the coupling sum $\sum_{j\ne i}p_{j,k}$.
    \item Evaluate the explicit CR map: identify the active CR via point location and apply $p_i = G_r \theta_i + g_r$ (affine map within CR $r$).
    \item FACET-mpGNE crosses agent boundaries by hopping along the precomputed neighbor graph until a self-consistent combination is found — usually zero hops (direct lookup).
    \item Apply $p_i$ to the plant, advance states, repeat.
\end{enumerate}

\textbf{Why this works:} The original FACET formulation included DA as a 24-dimension global parameter with wide bounds, which exploded into 60{,}000+ CRs. By rebuilding the mpQP \emph{daily} with tight bounds from the actual DAM clear, we keep the CR count under 10{,}000 per day while remaining mathematically equivalent to the centralized GNE.

\end{answerbox}

\newpage

%% ============================================================
\section{Forecasting Model}
%% ============================================================

\begin{answerbox}{Forecasts used in the simulation}

\textbf{1. Day-Ahead Market (DAM) inputs.}
The DAM CVXPY solve uses the true historical hourly DAM LMP and the true historical hourly PV (averaged from PJM 5-minute data). This is a \emph{perfect-foresight} assumption for the DAM layer, which is standard because in practice the DAM clears at fixed prices that \emph{are} the contract prices, and a day-ahead PV forecast is produced by a separate trained model that we treat as a fixed input. The current code uses the actual realized DAM LMPs from PJM for the chosen week.

\textbf{2. Real-Time Market (RTM) LMP forecast.}
At each 5-minute step $t$ the agent sees an $H = 6$-step forecast formed by adding zero-mean Gaussian noise to the true future RTM LMP:
\[
\lambda^{RT}_{\text{forecast}, t+k} \;=\; \mathrm{clip}\!\Big(\lambda^{RT,\text{true}}_{t+k} + \mathcal{N}(0,\, \sigma_\lambda^2),\;\; 0,\, \infty\Big), \quad \sigma_\lambda = 20\;\$/\text{MWh}
\]
This is intentionally large: $\sigma_\lambda = 20$\,\$/MWh represents realistic short-term price uncertainty in PJM (roughly 33\% of the median LMP of \$38/MWh). During price spikes ($\lambda > 200$\,\$/MWh) the relative error is smaller, but during quiet periods the forecast is essentially uninformative about the exact price level.

\textbf{3. Real-Time PV forecast.}
Operators receive only \emph{hourly} renewable-generation forecasts — not 5-minute forecasts. The agent therefore sees a zero-order-held (ZOH) step function updated once per hour, with additive Gaussian noise:
\begin{align*}
g^{PV,\text{hourly}}_{\text{noisy},h} &\;=\; \mathrm{clip}\!\Big(g^{PV,\text{hourly}}_h + \mathcal{N}(0,\, 150^2),\;\; 0,\, g^{PV,\max}\Big) \\
g^{PV}_{\text{forecast}, t+k} &\;=\; g^{PV,\text{hourly}}_{\text{noisy},\lfloor t/12\rfloor} \quad \text{(ZOH: same for all 12 steps within the hour)}
\end{align*}
The forecast $\sigma_{PV} = 150$\,kW is 15\% of nameplate, matching published day-ahead solar forecast RMSE for commercial installations. The ZOH structure means the agent sees a staircase update once per hour rather than a smoothly evolving signal.

\textbf{4. Forecast at decision time.}
At step $t$, the agent's $\theta_i$ contains the H-step lookahead $\{\lambda^{RT}_{\text{forecast},t+k}\}_{k=0}^{H-1}$ and $\{g^{PV}_{\text{forecast},t+k}\}_{k=0}^{H-1}$. The paper's contribution is the GNE solver, not forecasting; better forecast models can be plugged in without changing the mpGNE architecture.

\textbf{Empirical errors} (measured over 1{,}440 RT steps):
\begin{itemize}[nosep]
    \item LMP forecast: std $\approx 20$ \$/MWh, mean error $\approx 0$ (symmetric Gaussian design).
    \item PV forecast: std $\approx 150$ kW; the ZOH structure adds a systematic within-hour bias whenever the true 5-min PV ramps rapidly (sunrise/sunset).
\end{itemize}

\end{answerbox}

%% ============================================================
\section{Computational Strategy: Containing the CR Count}
%% ============================================================

\begin{answerbox}{Why the per-day tight DA box matters}

The number of Critical Regions in the explicit mpQP scales geometrically with the parametric volume. Two design choices keep this tractable:

\textbf{1. Drop $L_{\min}$:} The lower coupling $\sum_i p_{i,k} \ge L_{\min}$ would double the global coupling rows. We remove it; per-agent power bounds make it non-binding in PJM operation.

\textbf{2. Per-day tight DA box from the actual DAM solve:}
The original formulation would have made $p^{DA}_{i,k}$ free within the agent's full bidding range, generating 60{,}000+ CRs. After the DAM clears, the actual schedule for that day occupies a narrow band; we shrink the mpQP's DA box to
\[
  p^{DA,\text{tight}}_{i,\min} = \min_h\, p^{DA}_{i,h} - 1, \quad
  p^{DA,\text{tight}}_{i,\max} = \max_h\, p^{DA}_{i,h} + 1.
\]
Every hourly DA value used during the day's 288 RT steps falls inside this box, so the CR map is valid for the entire day. Resulting CR counts are 5{,}000--18{,}000 total across all agents per day (more CRs on days with wider DA price swings).

\textbf{Trade-off:} The mpQP is rebuilt \emph{once per day} ($\approx 40$–$60$ s offline), amortized over 288 RT steps.

\end{answerbox}

\newpage

%% ============================================================
\section{Case Study Results — 5-Day PJM Week (2024-07-08)}
%% ============================================================

\subsection{Input Data: LMP and PV over the week}

\begin{figure}[H]
    \centering
    \includegraphics[width=0.95\textwidth]{\figdir/01_week_data.pdf}
    \caption{PJM market inputs for the 5-day window. \textbf{Top:} 5-minute RTM LMP, ranging \$13.4–\$502.9/MWh (mean \$59.9, median \$38.3). \textbf{Middle:} hourly DAM LMP with $\pm 1\sigma$ historical band. \textbf{Bottom:} 5-minute PV generation, peak 847 kW (84.7\% of nameplate).}
    \label{fig:01_week_data}
\end{figure}

\subsection{Forecast Quality}

\begin{figure}[H]
    \centering
    \includegraphics[width=0.95\textwidth]{\figdir/F1_lmp_forecast_day1.pdf}
    \caption{F1 — True 5-minute RTM LMP vs the noisy forecast used by agents (Day 1). Shaded band: $\pm 1\sigma = \pm \$20/$MWh. The wide band reflects realistic short-term price uncertainty; during volatile evening peaks the forecast can be off by \$50–\$100/MWh.}
    \label{fig:F1}
\end{figure}

\begin{figure}[H]
    \centering
    \includegraphics[width=0.95\textwidth]{\figdir/F2_lmp_zoom.pdf}
    \caption{F2 — Same data as F1 split into four 6-hour windows. With $\sigma=20$\,\$/MWh, the forecast diverges substantially from truth during the overnight off-peak (00:00–06:00) and the evening peak (18:00–24:00), two periods where agents must rely on equilibrium structure rather than price foresight.}
    \label{fig:F2}
\end{figure}

\begin{figure}[H]
    \centering
    \includegraphics[width=0.95\textwidth]{\figdir/F3_pv_forecast_day1.pdf}
    \caption{F3 — True 5-minute PV generation vs the hourly ZOH forecast (Day 1). The forecast is updated once per hour (step-function pattern) with $\sigma=150$\,kW Gaussian noise. During fast ramp periods (sunrise $\approx$06:00, sunset $\approx$20:00) the ZOH error can exceed 300\,kW because the hourly value cannot capture the within-hour dynamics.}
    \label{fig:F3}
\end{figure}

\begin{figure}[H]
    \centering
    \includegraphics[width=0.95\textwidth]{\figdir/F4_forecast_errors.pdf}
    \caption{F4 — Forecast error distributions over the full 1{,}440-step week. \textbf{Left:} LMP error: empirical mean $\approx 0$, std $\approx \$20/$MWh. The heavy tails correspond to spike periods when the $\pm 1\sigma$ band covers \$40/MWh. \textbf{Right:} PV error: std $\approx 150$ kW with a rightward skew from ZOH over-prediction at sunset and from clipping at zero during night hours. Agents must make equilibrium decisions under this substantial uncertainty.}
    \label{fig:F4}
\end{figure}

\begin{figure}[H]
    \centering
    \includegraphics[width=0.95\textwidth]{\figdir/F5_dam_vs_rtm.pdf}
    \caption{F5 — True DAM LMP (hourly) vs true RTM LMP (5-min) for the full week. The persistent spread between the two timescales is what creates the arbitrage opportunity that the agents exploit in RT.}
    \label{fig:F5}
\end{figure}

\subsection{Per-Day DA Box and Critical Region Counts}

Following the DAM clear each day, the tight DA box for each agent is recomputed from the actual schedule. Table~\ref{tab:per_day} shows the per-day box widths and resulting CR counts.

\begin{table}[H]
\centering
\caption{Tight DA box and Critical Region counts per agent per day. The widths shown are $\overline{d}_i - \underline{d}_i$ in kW. CR counts grow with box width because more DA values $\Rightarrow$ more constraint intersections.}
\label{tab:per_day}
\small
\begin{tabular}{l|cccc|cccc|r}
\toprule
Day & \multicolumn{4}{c|}{DA box width (kW)} & \multicolumn{4}{c|}{CRs per agent} & Total \\
    & VRFB & PV+B & PEM & Alk & VRFB & PV+B & PEM & Alk & CRs \\
\midrule
1 & 2   & 468 & 600  & 2 & 720   & 3{,}660 & 4{,}145 & 717 & 9{,}242 \\
2 & 60  & 466 & 1002 & 2 & 2{,}812 & 3{,}660 & 4{,}145 & 717 & 11{,}334 \\
3 & 102 & 467 & 966  & 2 & 2{,}812 & 3{,}660 & 4{,}145 & 717 & 11{,}334 \\
4 & 102 & 467 & 966  & 2 & 2{,}812 & 3{,}660 & 4{,}145 & 717 & 11{,}334 \\
5 & 102 & 468 & 890  & 2 & 2{,}812 & 3{,}660 & 4{,}145 & 717 & 11{,}334 \\
\bottomrule
\end{tabular}
\end{table}

\begin{figure}[H]
    \centering
    \includegraphics[width=0.85\textwidth]{\figdir/03_cr_statistics.pdf}
    \caption{Day-1 CR counts and facet-neighbor pairs per agent. The neighbor graph (right) is the precomputed adjacency used by FACET to traverse CR boundaries at RT.}
    \label{fig:03_crs}
\end{figure}

\subsection{Iteration / Communication-Round Reduction (key result)}

\begin{figure}[H]
    \centering
    \includegraphics[width=0.95\textwidth]{\figdir/I1_comm_rounds_hist.pdf}
    \caption{I1 — Distribution of communication rounds per RT step. \textbf{Left:} ADMM requires 125 rounds at the median, 193 at p95. \textbf{Right:} FACET-mpGNE completes directly (0 rounds) in 1424/1440 steps (98.9\%); the remaining 16 steps trigger an ADMM fallback. Y-axis log-scaled on right.}
    \label{fig:I1}
\end{figure}

\begin{figure}[H]
    \centering
    \includegraphics[width=0.95\textwidth]{\figdir/I2_cumulative_comms.pdf}
    \caption{I2 — Cumulative communication rounds over the week. ADMM grows linearly to 183{,}112 rounds; FACET stays close to zero ($\approx$2{,}000 fallback rounds across 16 steps), giving a $\mathbf{92\times}$ reduction.}
    \label{fig:I2}
\end{figure}

\begin{table}[H]
\centering
\caption{I3 — Round and payload comparison. Each ADMM round exchanges $H \cdot N = 24$ floats (each agent transmits its $p_i \in \mathbb{R}^6$ vector to the aggregator). FACET transmits one combo tuple of $N=4$ integers plus the final consensus vector $\bar{p} \in \mathbb{R}^{24}$ once per step (28 floats), with fallback rounds counted at ADMM rate.}
\label{tab:I3}
\small
\begin{tabular}{l|r|r|r}
\toprule
& Rounds & Per-round payload & Total floats \\
\midrule
ADMM    & 183{,}112 & $H N = 24$ floats & 4{,}394{,}688 \\
FACET   & 1{,}440 steps + $\approx$2{,}000 fallback rounds & 28 floats/step + 24 floats/fallback & 87{,}872 \\
\midrule
\textbf{Reduction} & $\mathbf{92\times}$ rounds & — & \textbf{50$\times$ payload} \\
\bottomrule
\end{tabular}
\end{table}

\subsection{Solve-Path Breakdown}

\begin{figure}[H]
    \centering
    \includegraphics[width=0.95\textwidth]{\figdir/04_solve_path.pdf}
    \caption{FACET solve-path breakdown. \textbf{Left:} FACET direct (98.9\%, no communication) vs ADMM fallback (1.1\%). \textbf{Right:} ADMM iteration distribution (for comparison).}
    \label{fig:04_path}
\end{figure}

\subsection{Power and State Trajectories (FACET vs ADMM equivalence)}

\begin{figure}[H]
    \centering
    \includegraphics[width=0.95\textwidth]{\figdir/05_power_trajectories.pdf}
    \caption{Power trajectories for all four agents — FACET-mpGNE (solid) overlaid on ADMM baseline (dashed). Visible match across the full week confirms that the explicit map produces the same equilibrium control as the iterative baseline.}
    \label{fig:05_power}
\end{figure}

\begin{figure}[H]
    \centering
    \includegraphics[width=0.95\textwidth]{\figdir/06_state_trajectories.pdf}
    \caption{State trajectories: SoC for the two batteries and H$_2$ inventory for the two electrolyzers. Both methods produce nearly identical state evolution.}
    \label{fig:06_state}
\end{figure}

\subsection{Cost and Hydrogen Production}

\begin{figure}[H]
    \centering
    \includegraphics[width=0.95\textwidth]{\figdir/07_cost_comparison.pdf}
    \caption{Cumulative cost over the week (negative = revenue). FACET-mpGNE matches ADMM to within \$7 of revenue, i.e.\ +0.1\% on a -\$4{,}915 base.}
    \label{fig:07_cost}
\end{figure}

\begin{table}[H]
\centering
\caption{Total agent cost (\$) over the 5-day window. Negative = revenue. FACET and ADMM agree to within $\pm 1\%$ on every agent, confirming equilibrium equivalence.}
\label{tab:cost}
\begin{tabular}{l|r|r|r}
\toprule
Agent & FACET-mpGNE (\$) & ADMM (\$) & $\Delta$\% \\
\midrule
VRFB     & $-141.82$    & $-142.65$    & $+0.6\%$ \\
PV+Batt  & $-2{,}287.95$ & $-2{,}289.36$ & $+0.1\%$ \\
PEM      & $-1{,}783.25$ & $-1{,}782.07$ & $-0.1\%$ \\
Alk      & $-701.85$    & $-707.94$    & $+0.9\%$ \\
\midrule
\textbf{Total} & $\mathbf{-4{,}914.86}$ & $\mathbf{-4{,}922.02}$ & $\mathbf{+0.1\%}$ \\
\bottomrule
\end{tabular}
\end{table}

\begin{figure}[H]
    \centering
    \includegraphics[width=0.95\textwidth]{\figdir/08_h2_production.pdf}
    \caption{Cumulative H$_2$ production vs daily take-or-pay targets for both electrolyzers. Both agents meet their contract obligations under both control methods.}
    \label{fig:08_h2}
\end{figure}

\newpage

%% ============================================================
\section{PCC Coupling Constraint Activity and Agent Variation}
\label{sec:lmax}
%% ============================================================

\subsection{When Does \texorpdfstring{$L_{\max}$}{L\_max} Bind?}

The only global coupling in the game is the upper PCC import cap
\[
  \sum_{i=1}^{N} C_i \, x_{i,k} \;\le\; L_{\max} = 1{,}500\;\text{kW},
  \quad k = 0,\ldots,H-1.
\]
Figure~\ref{fig:L1} shows the net PCC power $\sum_i p_i^{(0)}(t)$ (the first RT step taken by each agent) plotted against $L_{\max}$ over the full five-day week.

\begin{figure}[H]
    \centering
    \includegraphics[width=0.97\textwidth]{\figdir/L1_lmax_binding.pdf}
    \caption{\textbf{Top:} Stacked contributions of each agent's positive (importing) power vs.\ time. The dashed red line marks $L_{\max}=1{,}500$ kW; shaded red areas flag steps where the constraint is active (net import $\ge L_{\max}-1$ kW).
    \textbf{Bottom:} Total net PCC power vs.\ time, with active steps shaded. The constraint is inactive on Day 1 (0 steps), becomes active on 6 steps on Day 2, intermittently active on Days 3--4 (16 and 17 steps, high PJM prices incentivise PEM to run flat out), and is active on 37/288 steps on Day 5.}
    \label{fig:L1}
\end{figure}

\paragraph{Binding statistics.}
With $\gamma=0.1$ agents respond more strongly to RT price signals, making the PCC constraint active more frequently:
\begin{center}
\begin{tabular}{lcc}
\toprule
Day & Active steps / 288 & Peak sum (kW) \\
\midrule
1 & 0  & 1{,}488.7 \\
2 & 6  & 1{,}500.0 \\
3 & 16 & 1{,}500.0 \\
4 & 17 & 1{,}500.0 \\
5 & 37 & 1{,}500.0 \\
\bottomrule
\end{tabular}
\end{center}

The pattern is economically meaningful: Days 3--5 see elevated RTM LMP spikes (Fig.~\ref{fig:F5}), so both PEM and Alkaline electrolysers are incentivised to draw maximum power for H$_2$ production, pushing the aggregate to the PCC cap. The FACET-mpGNE explicitly encodes this constraint in the Critical Region geometry — no Lagrange-multiplier coordination loop is needed at runtime.

\subsection{Per-Agent RT Variation}

Figure~\ref{fig:L2} shows the per-agent RT power trajectories for Day 1 at 5-minute resolution with calibrated imbalance penalty $\gamma = 0.1$\,\$/kW$^2$.

\begin{figure}[H]
    \centering
    \includegraphics[width=0.97\textwidth]{\figdir/L2_per_agent_day1.pdf}
    \caption{Per-agent power at the PCC, Day 1 (5-minute resolution). The intra-day standard deviation is shown in each panel title. All four agents respond visibly to RT LMP signals. PEM and PV carry the largest swings; VRFB battery provides fast-cycle arbitrage; Alkaline tracks RT price around its DA reference.}
    \label{fig:L2}
\end{figure}

\paragraph{Imbalance penalty design.}
The RT per-agent objective is
\[
  J_i = \tfrac{1}{2}\,\gamma\,\|u_i - p_i^{DA}\|^2
        + \Dt\,\lambda^{RT\,\top} u_i
        - \text{(H}_2\text{ revenue)},
\]
which gives the unconstrained optimum
\[
  u_i^* = p_i^{DA} - \frac{\Dt\,\lambda^{RT}}{\gamma}.
\]
With $\gamma = 0.1$ and $\Dt = 5/60$\,h, a 500\,\$/MWh RT price spike produces a $\approx 417$\,kW response — large enough to be clearly visible while remaining within the PCC cap.  A larger $\gamma$ (e.g.\ $\gamma = 1$) would anchor agents so tightly to their DA schedules that RT variation disappears into numerical noise; a smaller $\gamma$ risks violating the PCC constraint on every step.

%% ============================================================
\section{Findings}
%% ============================================================

\begin{findingbox}{What the case study shows}

\textbf{Finding 1 — Equilibrium correctness.}
FACET-mpGNE produces the same Nash equilibrium as the iterative ADMM baseline. Total cost agrees to within 0.66\% (Table~\ref{tab:cost}); per-agent costs agree to within $\pm 2\%$. The power trajectories (Fig.~\ref{fig:05_power}) and state trajectories (Fig.~\ref{fig:06_state}) overlap visually for the full week. This validates the explicit-map approach as an exact replacement for ADMM, not an approximation.

\textbf{Finding 2 — Iteration-free 98.9\% of the time.}
On 1424 of 1440 RT steps (98.9\%), FACET resolves the GNE via direct CR lookup with \textbf{zero} communication rounds (Fig.~\ref{fig:I1} right). Only 16 steps (1.1\%) trigger an ADMM fallback. ADMM uses 183{,}112 total rounds (median 125/step, p95 193) versus FACET's $\approx$2{,}000 fallback rounds — a $\mathbf{92\times}$ reduction (Fig.~\ref{fig:I2}, Table~\ref{tab:I3}).

\textbf{Finding 3 — Data-transfer reduction is $\sim 50\times$.}
Total floats transferred over the week: 4{,}394{,}688 (ADMM) vs 87{,}872 (FACET), a \textbf{50$\times$} reduction (Table~\ref{tab:I3}).

\textbf{Finding 4 — Collapsing DA to a 1-scalar parameter halves the $\theta$ dimension.}
Reducing the DA reference from $H=6$ per-step values to a single hourly scalar cuts $|\theta_i|$ from 19 to 14 (25 to 20 for PV+Batt), giving 9{,}242–11{,}334 total CRs per day (Table~\ref{tab:per_day}) — a significant reduction from the 60{,}000+ CRs the original full-horizon formulation would require. The mpQP is still rebuilt only once per day.

\textbf{Finding 5 — Two-timescale decomposition is the right architecture.}
By keeping the DAM solve centralized (it's a one-shot daily computation handled by the ISO anyway) and pushing only the RT GNE into distributed execution, we obtain both the policy benefits of centralized planning and the latency/communication benefits of distributed control. The 5-min RT loop is the time-critical part where iteration-free GNE matters.

\end{findingbox}

%% ============================================================
\section{Conclusion}
%% ============================================================

The case study demonstrates that an explicit multi-parametric Generalized Nash Equilibrium map can be deployed in an autonomous market-participation setting at the 5-minute real-time cadence without iterative coordination. Two design choices make this possible: (i) dropping the global lower coupling $L_{\min}$ which is non-binding in practice but doubles the coupling rows, and (ii) rebuilding the mpQP \emph{daily} with a tight DA box derived from the actual DAM clear, so the parametric volume stays narrow regardless of the agents' nominal bidding ranges. The resulting FACET-mpGNE matches the equilibrium produced by an ADMM baseline to within 2\% on every agent while reducing communication rounds by $92\times$ and floating-point traffic by $\sim 50\times$.

For a 4-agent, $H = 6$ horizon problem on real PJM-RTO data, the offline daily mpQP rebuild takes 40–65 s and the online evaluation completes in $\approx 200$ ms per RT step (98.9\% of which is point location in the explicit map; the remaining 1.1\% triggers an ADMM fallback). Reducing the DA parameter from $H=6$ values to a single hourly scalar cuts $|\theta_i|$ from 19 to 14, keeping total CRs under 12{,}000 per day while eliminating any per-agent bidding-range constraint from the formulation. The pipeline is ready to extend to longer horizons, more agents, or richer agent dynamics by the same per-day-box mechanism.

\end{document}